\documentclass[UTF8]{article}
\usepackage[a4paper,margin=1in]{geometry}
\usepackage{booktabs}
\usepackage{hyperref}
\usepackage{listings}
\usepackage{graphicx}
\usepackage{ctex}

\title{基于CUDA + pybind11的简易卷积网络框架与CIFAR-10分类}
\author{2400013080 李玮博}
\date{\today}

\begin{document}
\maketitle

\section{实验背景}
深度学习框架通常将张量运算、自动微分与GPU加速高度封装。为了理解卷积网络的底层实现流程，
本实验在不依赖现成深度学习框架的前提下，使用CUDA实现卷积前向与反向传播，
通过pybind11暴露为Python可调用模块，并在Python侧完成自动微分与优化器。
最终在CIFAR-10数据集上完成图像分类任务，验证该最小框架能够完成端到端训练。

\section{方法}
\subsection{整体架构}
框架分为三层：
（1）CUDA/C++底层：实现张量存储、基础算子、卷积与反向传播等；
（2）pybind11绑定层：将CUDA算子注册为Python可调用接口；
（3）Python前端：构建Tensor与计算图，完成自动微分、优化器与模型定义。

\subsection{CUDA卷积前向与反向}
卷积实现位于 \texttt{csrc/layers/conv.cu}，采用 \texttt{im2col} 将输入展开为矩阵，
再调用GPU矩阵乘法完成卷积前向。反向传播包含两部分：
对卷积核的梯度通过 \texttt{col} 与输出梯度矩阵相乘得到；
对输入的梯度通过转置卷积（\texttt{col2img}）重建并累加。
该实现支持 padding 与 stride。

\subsection{pybind11接口}
绑定代码位于 \texttt{csrc/pybind.cpp}，编译为 \texttt{torchofdreaveler.core} 模块。
接口以函数形式暴露，例如 \texttt{convolve} 与 \texttt{convolve\_backward}，并在Python侧封装为算子类。

\subsection{自动微分与优化器}
Python侧 \texttt{torchofdreaveler/\_core/operators.py} 定义计算图与算子，
通过 \texttt{compute\_gradient\_of\_variables} 反向传播梯度。
优化器实现位于 \texttt{torchofdreaveler/optim}，提供SGD与Adam。
损失函数采用Softmax交叉熵，CUDA实现位于 \texttt{csrc/layers/softmax\_loss.cu}。

\subsection{模型与训练流程}
模型示例采用Lenet风格的CNN与ResNetSmall（见 \texttt{models/lenet.py} 与 \texttt{models/resnet.py}），
训练脚本为 \texttt{train.py}。
数据集加载通过 \texttt{torchofdreaveler/data} 封装 torchvision 的 CIFAR-10 下载与预处理。

\section{实验结果}
表~\ref{tab:results} 给出一组可复现实验结果（来源：\texttt{result.md} 中的训练记录）。
ResNet20 在 CIFAR-10 上达到约 89.5\% 测试准确率；简单CNN基线约 75\%。

\begin{table}[h]
\centering
\caption{CIFAR-10 分类结果}
\label{tab:results}
\begin{tabular}{lccc}
\toprule
模型 & 训练轮数 & 最佳测试准确率 & 备注 \\
\midrule
ResNet20 & 100 & 89.5\% & 自定义框架训练 \\
Simple CNN & 40 & 75.1\% & 参考基线 \\
\bottomrule
\end{tabular}
\end{table}

\begin{figure}[h]
\centering
\includegraphics[width=0.85\linewidth]{resnet20_acc.png}
\caption{ResNet20 CIFAR-10 accuracy curve}
\label{fig:resnet20_acc}
\end{figure}

\begin{figure}[h]
\centering
\includegraphics[width=0.85\linewidth]{simplecnn_acc.png}
\caption{Simple CNN CIFAR-10 accuracy curve}
\label{fig:simplecnn_acc}
\end{figure}

\begin{figure}[h]
\centering
\includegraphics[width=0.85\linewidth]{simplecnn_loss.png}
\caption{Simple CNN CIFAR-10 loss curve}
\label{fig:simplecnn_loss}
\end{figure}


\section{讨论}
（1）\textbf{正确性验证}：卷积前向与反向与PyTorch对齐测试已在 \texttt{tests/test\_operator\_parity.py} 中实现，
为算子正确性提供保障。 \\
（2）\textbf{性能瓶颈}：当前实现以 \texttt{im2col + GEMM} 为主，内存开销较大，
且未对数据布局与融合算子做进一步优化。 \\
（3）\textbf{模型规模与精度}：简单网络收敛快但精度上限有限，
ResNet20 能较好展示框架对更复杂网络的支持能力。

\section{结论}
本实验完成了基于CUDA的卷积前向/反向、pybind11绑定、Python自动微分与优化器，
构建了一个可训练的轻量级卷积网络框架，并在CIFAR-10上验证了可用性。
后续可进一步优化算子性能、扩展算子覆盖率并完善工程化工具链。

\section{如何运行与复现}
\subsection{环境准备}
\begin{itemize}
  \item Windows + CUDA Toolkit（示例使用 CUDA 13.x）
  \item Python 3.8+
  \item 依赖：\texttt{torch}, \texttt{torchvision}, \texttt{numpy}
\end{itemize}

\subsection{pip安装包}
可以直接通过pip安装
\begin{lstlisting}[language=bash]
pip install .
\end{lstlisting}

\subsection{包的使用}
与torch的语言类似，具体细节可以直接参考\texttt{train.py}文件以及所有train开始的文件。

且当前实现支持自定义sgd adam等优化器以及scheduler的参数 并支持checkpoint保存模型，支持使用tensorboard保存记录并进行可视化
\subsection{训练CIFAR-10}
\begin{lstlisting}[language=bash]
python train.py --model lenet --epochs 50 --batch-size 512 --lr 0.1
\end{lstlisting}
该命令会自动下载 CIFAR-10 至 \texttt{./datasets/cifar10}，并在 \texttt{checkpoints/} 输出模型权重。
若使用ResNet20复现实验结果，可运行：
\begin{lstlisting}[language=bash]
python train.py --model resnet20 --epochs 100 --batch-size 512 --lr 0.1
\end{lstlisting}

\section{bonus：训练ImageNet}
由于2012版下载包过大，因此只使用Tiny ImageNet来进行训练，我尝试实现了BatchNorm，但实现细节可能出现问题，
算法并行度不够高，导致加上BN后GPU的使用率极低，而不加BN的resnet18在学习率较大时极易爆炸，经常出现loss为nan的情况，
因此仍然使用resnet20去实现，但由于resnet20的参数量不够使得不太能拟合好模型，且同样由于缺少BN，收敛速度较慢。
结果为训练了136个epoch后测试准确率为38.14\% 

\begin{figure}[h]
\centering
\includegraphics[width=0.85\linewidth]{resnet20_ti_acc.png}
\caption{ResNet20 Tiny ImageNet accuracy curve}
\label{fig:resnet20_ti_acc}
\end{figure}

\begin{figure}[h]
\centering
\includegraphics[width=0.85\linewidth]{resnet20_ti_loss.png}
\caption{ResNet20 Tiny ImageNet loss curve}
\label{fig:resnet20_ti_loss}
\end{figure}


具体训练模型与训练参数细节可以参照\texttt{train\_tinyimagenet\_renet20.py}

\end{document}
